\documentclass[10pt,journal]{IEEEtran}

\usepackage{cite}
\usepackage[cmex10]{amsmath}
\usepackage{amssymb,amsfonts}
\usepackage{booktabs}
\usepackage{multirow}
\usepackage{graphicx}
\usepackage{url}
\usepackage[hidelinks]{hyperref}

\newcommand{\Method}{SFC-VSR}
\newcommand{\Dataset}{Chinese-LiPS}

\begin{document}

\title{Sparse-Feedback Continual Personalization for Mandarin Visual Speech Recognition: Gains, Feedback Corruption, and the Limits of Dynamic Experts}

\author{Zezhou~Wang, Jun~Li, Xiaoming~Li, and Congxin~Wei%
\thanks{Z. Wang and C. Wei are with the School of Computer Science and Artificial Intelligence, Lanzhou University of Technology, Lanzhou, China.}%
\thanks{J. Li is with the College of Computer Science and Engineering, Northwest Normal University, Lanzhou, China.}%
\thanks{X. Li is with the Department of Otolaryngology Head and Neck Surgery, 980th Hospital of Logistic Supportive Force of Chinese PLA, Shijiazhuang, China.}%
\thanks{Corresponding author: Congxin Wei.}}

\maketitle
\bstctlcite{IEEEexample:BSTcontrol}

\begin{abstract}
Visual speech recognition (VSR) models are usually evaluated as static predictors, although deployed users and conditions vary over time. We study continual personalization of a frozen Mandarin VSR backbone when one character-level correction is available every ten utterances. Our strictly prequential protocol scores each transcript before any correction can update a 75k-parameter residual adapter; reliability checks admit pseudo targets and post-update checks permit atomic rollback. On a 3,908-utterance Chinese-LiPS test stream, no-feedback adaptation lowers the three-seed mean character error rate (CER) from 69.30\% to $67.71\pm0.18$\%. Sparse corrections further reduce it to $66.88\pm0.10$\%, a 0.832-point mean improvement with the same direction at all three matched seeds. Dynamic experts improve over one adapter by only 0.055 points and are retained as pre-fix diagnostics because feedback location could affect growth before prediction. Error localization and active querying also fail pre-specified development gates, leaving the frozen holdout unread. A five-arm, single-seed audit on a separate 700-utterance development stream then tests stronger matched controls. BN-TENT-VSR and source-free ETA-VSR worsen static CER by 36.755 and 38.876 points. Under the same 70-correction budget and similar parameter counts, combined replay is 1.229 points better than online LoRA (paired 95\% CI: 0.737 to 1.720) and improves static-corrected revisit forgetting by 2.290 points (0.599 to 3.970). This last comparison is budget- and size-matched, not update-source matched, because replay also accepts 11 pseudo updates. The evidence supports sparse-feedback replay as the development incumbent, but not the present expert, localized, active-query, normalization-only, online-LoRA, or Feature-FiLM alternatives; the development-only audits are not confirmatory holdout results.
\end{abstract}

\begin{IEEEkeywords}
visual speech recognition, lip reading, continual adaptation, personalization, sparse feedback, test-time adaptation
\end{IEEEkeywords}

\section{Introduction}
\label{sec:introduction}

\IEEEPARstart{V}{isual} speech recognition converts silent mouth motion into text and enables communication when audio is unavailable, private, or unreliable~\cite{potamianos2003recent,ma2022visual}. Modern systems combine a spatiotemporal visual front end with a sequence encoder and CTC~\cite{graves2006ctc} or attention-based decoding. Their reported accuracy, however, normally assumes a fixed model. A deployed recognizer instead observes a stream: speakers recur, camera conditions drift, and a user may occasionally correct a transcript. The practical question is therefore not only whether a source model generalizes, but whether a small amount of feedback can improve future predictions without replaying source data or corrupting previously useful behavior.

Speaker adaptation has clear precedent in VSR. Prompt- and LoRA-based personalization can adapt both visual and language components to a target user~\cite{yeo2024personalized}. Continual test-time adaptation (CTTA) similarly studies non-stationary unlabeled streams, including parameter-efficient adaptation, recovery mechanisms, and expert routing~\cite{niu2022eata,wang2022cotta,lee2024becotta,cui2025paint}. These settings do not directly answer our question. VSR personalization is commonly organized as an offline support-set stage, while generic CTTA is primarily evaluated on image classification or semantic segmentation. Sparse human corrections introduce a different resource: labels are direct but intermittent and may themselves be noisy, and every result must be recorded before its label is used.

This paper studies that setting with a frozen Mandarin VSR backbone and lightweight residual adapters. The strict causal path uses one adapter and separates prediction from any later correction. We also audit an earlier dynamic-expert prototype under matched streams, but its known feedback-location flag could affect pre-prediction growth; all expert numbers are therefore retained only as pre-fix diagnostics. This distinction matters because a more complex structure is useful only if it routes domains meaningfully and improves either accuracy or revisit forgetting beyond a single adapter.

Our contributions are empirical and methodological:
\begin{enumerate}
  \item We define an auditable sparse-feedback protocol for continual Mandarin VSR. Predictions are scored before adaptation, stream order and feedback positions are fixed, manifests and checkpoints are hash-locked, and comparisons use utterance-paired bootstrap intervals.
  \item We separate two measured gains: no-feedback online adapter adaptation improves seed-42 CER by 1.72 points, while one correction every ten utterances produces a stable three-seed total gain of 2.42 points over the static model.
  \item We provide a negative diagnostic for the current expert implementation under two routing regimes: the original expert bank collapses almost entirely to one expert, while a validation-calibrated repair over-fragments into eight experts and performs materially worse than one adapter. These pre-fix runs do not establish that dynamic experts are intrinsically ineffective.
  \item We audit where correction gradients should act. Greedy correct-span preservation is worse than whole-sequence replay. Target-conditioned CTC localization significantly beats same-mass randomized support, yet remains significantly worse than replay and is rejected by its validation gate.
\end{enumerate}

The paper deliberately separates supported claims from open work. The validation-only router gate is resolved as a no-go, so no repaired-router test run or additional expert seed is reported. An adapter-only entropy baseline is included as a failure diagnostic; stronger matched CTTA baselines remain required before a method-comparison submission.

\section{Related Work}
\label{sec:related_work}

\subsection{Mandarin Visual Speech Recognition}

Large-scale VSR systems use 3D convolutional visual front ends and Transformer or Conformer sequence models~\cite{ma2022visual}. Chinese VSR adds a large character vocabulary and severe visual ambiguity. Chinese-LiPS provides 100.84 hours from 207 speakers, including fixed validation and test splits of 1,959 and 3,908 clips~\cite{zhao2025chineselips}. We use its lip video and transcript channels only. The task is cross-dataset deployment: the source checkpoint is fixed before Chinese-LiPS adaptation, and no Chinese-LiPS training sample is used by the online learner.

\subsection{VSR Personalization}

Personalized Lip Reading adapts a pretrained recognizer at visual and language levels using prompt tuning and LoRA, and evaluates on the speaker-adaptation dataset VoxLRS-SA~\cite{yeo2024personalized}. That work establishes the value of parameter-efficient speaker adaptation. Our setting differs in evaluation time and supervision: corrections arrive inside a non-stationary stream, every prediction is evaluated before updating, and returned domains test whether prior adaptation survives intervening speakers.

\subsection{Continual Test-Time Adaptation}

EATA filters unreliable or redundant samples and regularizes adaptation to reduce forgetting~\cite{niu2022eata}; CoTTA uses augmentation, teacher averaging, and stochastic restoration for continual shifts~\cite{wang2022cotta}. AETTA estimates online accuracy without labels and uses the estimate to control recovery~\cite{lee2024aetta}. These methods motivate reliability and recovery controls; the cited mechanisms do not by themselves establish novelty for per-update rollback.

Expert-based CTTA explicitly addresses domain interference. BECoTTA blends low-rank domain experts with input-dependent routing~\cite{lee2024becotta}. PAINT uses a query mechanism to select a known domain prompt or allocate a new one~\cite{cui2025paint}. Consequently, adapter banks, cosine routing, and dynamic allocation cannot be claimed as first contributions here. Our expert bank is an evaluated VSR instantiation whose value must be established against one adapter and, in future work, against matched implementations of these nearest neighbors.

\section{Method}
\label{sec:method}

\subsection{Prequential Sparse-Feedback Setting}

Let $\{(x_t,y_t,d_t)\}_{t=1}^{T}$ be a video stream, where $d_t$ is retained only for ordering and post-hoc analysis. At step $t$, the system first predicts $\hat y_t$ and records its edit distance to $y_t$. Only then may it receive feedback $y_t$ according to a schedule fixed before evaluation. The main operating point supplies one correction every ten utterances. Character error rate is aggregated as
\begin{equation}
  \mathrm{CER}=\frac{\sum_t \operatorname{EditDistance}(\hat y_t,y_t)}{\sum_t |y_t|}.
\end{equation}
This ordering prevents the current correction from improving the score already assigned to the same sample.

\subsection{Frozen Backbone and Residual Adapter}

The base VSR model contains a visual encoder $f_{\theta}$ and CTC/attention decoder $h_{\theta}$; all base parameters $\theta$ remain frozen. For encoded features $z_t=f_{\theta}(x_t)$, a residual bottleneck adapter is
\begin{equation}
  A_{\phi}(z)=z+sW_{\mathrm{up}}\,\sigma\!\left(W_{\mathrm{down}}\operatorname{LN}(z)\right),
\end{equation}
where the bottleneck width is 48, the feature width is 768, and $s$ is a learnable scalar initialized to one. Zero initialization of $W_{\mathrm{up}}$ makes a newly created adapter function-preserving at creation. It does not guarantee preservation after subsequent updates.

When feedback is present, its (possibly corrupted) token sequence is used directly as the adaptation target and bypasses the pseudo-label admission gate. Without feedback, a pseudo target is admitted only if the CTC confidence, perturbed-view agreement, CTC emission rate, and CTC/attention agreement pass fixed reliability tests. One AdamW step minimizes a weighted combination of CTC target loss, clean/perturbed posterior consistency, entropy, and feature anchoring. The learning rate is $5\times10^{-4}$ and gradient norm is clipped at 1.

\subsection{Transactional Acceptance}

Before an update, the adapter and optimizer states are copied. A candidate update is rejected when target loss increases beyond a fixed tolerance or KL divergence from the frozen prediction exceeds its bound. When reliability checks are enabled, a pre-accepted candidate can also be rejected if post-update reliability fails, and a pseudo-target candidate can be rejected if its reliability score drops beyond the fixed tolerance. Rejection restores both parameter and optimizer states. This mechanism provides an atomic implementation property; whether it predicts future CER harm is an empirical question evaluated by ablation.

\subsection{Dynamic-Expert Validation}

The expert variant forms a family of signatures by concatenating frozen-feature mean and standard deviation with up to first- and second-order temporal-difference magnitudes. The normalized signature is routed by cosine similarity to an online prototype. Without feedback, low-similarity sequences are accumulated before growth; a feedback item can confirm growth immediately. Each accepted update also updates the selected prototype. For the single validation repair, order-2 signatures are extracted once from the frozen backbone and orders 0--2 are compared using a fixed leave-own-domain-out proxy. This selects order 1 and its cosine threshold without reading test outcomes. The proxy does not estimate true unseen-domain performance.

We test two necessary conditions: (i) expert routing must avoid collapse and reuse the first speaker's expert on return; and (ii) its overall CER or static-corrected revisit forgetting must materially improve over one adapter. Failure of either condition stops further dynamic-expert evaluation under the pre-specified decision rule.

\section{Experimental Protocol}
\label{sec:experiments}

\subsection{Data and Backbone}

Chinese-LiPS contains 36,208 clips from 207 speakers~\cite{zhao2025chineselips}. We use the official 3,908-clip test split (21 speakers) as an exploratory benchmark. The 1,959-clip validation split (11 disjoint speakers) was subsequently used for update-rule and routing selection and is therefore no longer an untouched confirmation set. The existing test matrix likewise includes feedback-cadence and safety explorations. Any next-round candidate is first subject to an explicit UID- and speaker-overlap audit and is evaluated on newly provenance-locked development and holdout streams; no result is promoted from development to test without a fixed holdout decision. Final external confirmation still requires a second, previously unused real shift stream. The source model is the public CNVSRC VSR checkpoint trained on CN-CVS, CN-CVS2/3, and CNVSRC. It uses a 3D CNN/ResNet visual front end, a 12-layer Conformer encoder, and CTC/bidirectional-attention decoding~\cite{cnvsrc2025baseline}.

Chinese-LiPS token IDs do not share the source vocabulary. We therefore re-encode original transcripts with the fixed source character vocabulary. Characters outside that vocabulary are dropped before any experiment; the test rate is 8.794\% (14,460 of 164,431 raw characters). The source CSV, text metadata, target vocabulary, generated manifest, and base checkpoint are SHA-256 locked. This normalization changes the evaluated target alphabet and is a limitation, not a performance optimization.

For the next-round candidate gates, we repurpose nine speakers from the official Chinese-LiPS training split that were never used by the preceding RSP-VSR experiments. The publicly documented source-checkpoint training inventory lists CN-CVS, CN-CVS2/3, and CNVSRC and does not list Chinese-LiPS; that inventory is not expanded into a sample-level absolute non-overlap proof. This is a protocol change, not an unused portion of the official validation set. Target-dev2 uses speakers 015/098/133 in an A--B--C--A stream of 154/253/243/154 utterances (804 total). Target-holdout2 uses disjoint speakers 046/001/093 and 136/242/264/136 utterances (778 total). Target-dev3, used only for the active-query gate and update-source completion, uses speakers 120/176/183 and 135/247/242/134 utterances (758 total). These pools have zero UID and speaker overlap with one another and with the historical validation and test streams. Their source CSVs, manifests, sidecars, vocabulary, and checkpoint are hash-locked before any candidate run; holdout2 remains unread by result analysis until a development rule has selected and frozen one candidate.

\subsection{Streams and Controls}

The main stream groups clips by speaker, shuffles speaker order and within-speaker order using seed 42, and never exposes the speaker label to the router. Historical feedback locations are fixed at every 10, 20, 50, or 100 samples. Main comparisons use seeds 7, 42, and 123 with the same sample order. In the corruption condition, each feedback character is independently replaced with probability 0.1 using the fixed run seed. On target-dev3, periodic, random, and uncertainty policies each receive exactly 75 corrections across the 75 complete ten-sample windows; the final eight samples receive no correction. Static, pseudo-only, feedback-only-periodic, and combined-periodic controls complete the same-stream update-source comparison.

The test revisit stream is A--B--C--A with 133, 255, 222, and 132 clips. The validation revisit uses three validation speakers and lengths 130, 195, 186, and 130. All UIDs are unique, so A1 and A2 contain different utterances from the same speaker. We report raw $\mathrm{CER}_{A2}-\mathrm{CER}_{A1}$ and subtract the corresponding static difference to control for sample difficulty. Method comparisons use paired utterance bootstrap for overall CER. Revisit difference-in-differences independently resamples paired method outcomes within A1 and A2.

\subsection{Compared Systems}

\begin{itemize}
  \item \textbf{Static}: frozen source model; no update.
  \item \textbf{Pseudo-only adapter}: one shared residual adapter with reliability-gated unlabeled updates and no corrections (required with the static, feedback-only, and combined arms for a complete update-source split; may be pending on a given stream).
  \item \textbf{Feedback-only adapter}: the same adapter updated only by scheduled ground-truth corrections; non-feedback samples are never used for parameter updates.
  \item \textbf{Combined replay}: reliability-gated pseudo updates plus full-sequence replay at the same correction locations. Pairing this arm with feedback-only alone yields a partial update-source comparison; same-stream pseudo-only is still needed for a full causal decomposition.
  \item \textbf{Expert bank$^{\dagger}$}: prototype routing, quarantine, and at most eight adapters.
  \item \textbf{Greedy correct-span}: full-sequence correction with KL preservation on greedily matched token segments.
  \item \textbf{Error-localized correction}: target-conditioned CTC error occupancy, with a same-mass randomized-support control.
  \item \textbf{Full-replay error hybrid}: combined replay plus edit-conditioned occupancy as an auxiliary term, with a matched randomized-support hybrid.
  \item \textbf{Adapter-only TENT}: entropy minimization~\cite{wang2021tent} on either all CTC frames or only nonblank greedy frames, updating the same adapter while the backbone remains frozen.
  \item \textbf{Safety ablations}: reliability disabled or rollback disabled.
\end{itemize}

$^{\dagger}$All expert-bank runs are pre-fix diagnostics in which the known feedback-location flag could influence growth before prediction; they are not strict causal comparisons.

The adapter-only TENT implementation is a matched diagnostic, not a reproduction of full-model TENT. Stronger TTA and expert baselines such as EATA, CoTTA, BECoTTA, and PAINT are not yet implemented on this VSR backbone. They remain required before a final method-comparison submission and are not replaced by numbers quoted from other datasets.

\subsection{Statistical Decision Rules}

Pairwise comparisons use 10,000-sample paired bootstrap 95\% intervals, and three-seed results include the population standard deviation. The error-localized candidate may enter test evaluation only if it improves validation CER over full-sequence replay by at least 0.3 points with a paired interval wholly below zero and is also significantly better than its same-mass randomized-support control. That gate has failed and is not reopened on the consumed validation stream. The full-replay hybrid is a new candidate: on the locked development stream it must improve replay by at least 0.3 CER points with a paired interval wholly below zero and must significantly beat its randomized-support hybrid. Its objective and weights are then frozen; the same direction, magnitude, and interval criteria must hold on the locked holdout before any test run or seed expansion. Feedback-only versus combined replay is reported as a partial update-source comparison on a shared correction schedule, not as a complete pseudo/feedback/interaction decomposition and not as a comparison across unequal annotation budgets. Same-stream static and pseudo-only arms remain required before any claim of a complete factorial split.

The active-query candidate is selected only on target-dev3. Uncertainty must improve combined-periodic CER by at least 0.3 percentage points with a paired 95\% interval wholly below zero, significantly outperform combined-random, and query a higher true-error rate than random with the paired-window interval wholly above zero. Each policy must spend exactly 75 corrections, and uncertainty minus periodic static-corrected A2--A1 forgetting must have a paired interval whose upper bound is at most zero. Failure of any condition stops the route without inspecting holdout2, adding seeds, or tuning the query rule. Passing freezes uncertainty as the only candidate allowed to accompany a fixed periodic baseline on holdout2; holdout outcomes cannot be used for further selection.

The dynamic-expert route proceeds beyond validation only if validation-only signature calibration reaches at least 0.8 proxy AUROC, F1, and clustering purity; the revisit router must use expert~0 for less than 90\% of samples and reuse the returning-A expert for at least 80\%; and expert-vs-single revisit improvement must be at least 0.3 CER percentage points in the better direction with a paired confidence interval whose upper bound is below zero. The calibration scores are leave-own-domain-out proxies, not evidence of true unseen-domain detection.

\section{Results and Analysis}
\label{sec:results}

\subsection{Online Adaptation and Sparse Feedback Both Improve CER}

Table~\ref{tab:main} summarizes the primary operating points on the complete test stream. At seed 42, no-feedback online adaptation improves static CER by 1.718 points. Adding feedback every ten samples improves the same-seed single adapter by a further 0.831 points and yields a 2.421-point three-seed mean gain over static. The three-seed standard deviation is 0.105 points, indicating that the total gain is not driven by one seed. More frequent corrections also matter: the expert-bank mean improves by 0.490 points from feedback every 20 to every 10 samples.

\begin{table}[t]
\centering
\caption{Exploratory CER on the 3,908-clip test stream. Mean $\pm$ population standard deviation is shown for three-seed rows; $\dagger$ marks pre-fix expert diagnostics.}
\label{tab:main}
\begin{tabular}{lcc}
\toprule
Method & Feedback & CER (\%) $\downarrow$ \\
\midrule
Static & none & 69.302 \\
Single adapter & none & 67.584 \\
Expert bank$^{\dagger}$ & none & 67.584 \\
Single adapter & every 20 & $67.335\pm0.063$ \\
Expert bank$^{\dagger}$ & every 20 & $67.317\pm0.058$ \\
Single adapter & every 10 & $66.881\pm0.105$ \\
Expert bank$^{\dagger}$ & every 10 & $66.826\pm0.079$ \\
\bottomrule
\end{tabular}
\end{table}

The no-feedback single and expert systems are identical at seed 42 because neither run grows a second effective expert. Their 1.718-point improvement over static is nevertheless a measured U0 result, although it currently has only one seed and lacks matched TTA baselines. The additional feedback gain belongs to a separate F10 supervision track and is supported by three seeds. We therefore report both effects without treating their different supervision budgets as interchangeable.

\subsection{Two Intuitive Update Alternatives Fail on Validation}

Full-sequence feedback replay reaches 69.637\% CER on the 641-utterance validation revisit, improving the 72.329\% static result by 2.692 points (paired 95\% CI: $-3.315$ to $-2.070$). Preserving greedily matched spans appears more selective, but reaches 70.174\% and is 0.538 points worse than replay (95\% CI: $+0.074$ to $+1.077$). The greedy correct-span variant therefore fails its validation gate and is not evaluated on test.

Adapter-only entropy minimization is substantially less stable. Selecting all CTC frames produces 98.581\% CER, while selecting only nonblank greedy frames produces 96.742\%. Relative to static, the two variants degrade CER by 26.252 and 24.413 points, respectively, and both paired intervals lie wholly in the harmful direction. We stop this baseline after validation rather than tune it against test outcomes. These failures do not imply that all TTA methods fail for VSR; they show that direct entropy minimization of this frozen-backbone adapter is not a competitive control.

\subsection{Error Localization Is Informative but Underperforms Replay}

Table~\ref{tab:validation_local} resolves the final validation-only update gate. Error localization reaches 72.004\% CER. It is 2.508 points better than randomized support (95\% CI: $-3.087$ to $-1.955$), showing that where the sparse correction gradient acts contains useful information. However, it is 2.367 points worse than full-sequence replay (95\% CI: $+1.757$ to $+3.048$). Its revisit difference-in-differences is also 3.244 points worse than replay (95\% CI: $+1.527$ to $+5.044$). Relative to static, the 0.325-point CER improvement is not significant (95\% CI: $-0.844$ to $+0.227$). The pre-specified replay gate therefore fails and no localized test run, parameter adjustment, or additional seed is conducted.

\begin{table}[t]
\centering
\caption{Validation correction-localization results. Corrected forgetting subtracts the static A2--A1 difference; lower is better.}
\label{tab:validation_local}
\begin{tabular}{lrrrr}
\toprule
Method & Overall & A1 & A2 & Corrected \\
 & \multicolumn{4}{c}{CER (\%)} \\
\midrule
Static & 72.329 & 58.505 & 62.713 & 0.000 \\
Full replay & 69.637 & 57.250 & 56.712 & $-4.746$ \\
Error-localized & 72.004 & 58.386 & 61.092 & $-1.502$ \\
Random support & 74.512 & 58.475 & 63.623 & $+0.940$ \\
\bottomrule
\end{tabular}
\end{table}

All 64 scheduled corrections receive a target-conditioned occupancy. Across the localized run, they contain 1,199 substitution and 143 deletion target positions plus 51 insertion frames, or 0.367\% of the 13,889 CTC frames. Mean matched and error occupancy masses are 11.204 and 21.007, with effective supports of 11.681 and 22.625 frames. The local objective decreases on all 64 corrections by 0.112 on average, yet total updates are 71 accepted, 3 rolled back, and 567 skipped. The randomized trajectory accepts only 27 updates and rolls back 46. Thus local objective descent is not sufficient evidence of future transcript improvement, and whole-sequence replay remains the selected update rule.

\subsection{Dynamic Experts Do Not Beat One Adapter Materially}

At feedback every ten samples, the expert advantage over one adapter is 0.087, 0.025, and 0.053 CER points for seeds 7, 42, and 123, respectively, for a mean of only 0.055 points. This is below the 0.3-point materiality threshold used for the structural decision. Because these prototype runs allowed the known feedback-location flag to confirm growth before the corresponding prediction, we do not treat the expert rows as strict causal comparisons; the negative routing and accuracy outcomes are diagnostic despite that favorable information.

The revisit stream leads to the same conclusion (Table~\ref{tab:revisit}). Both adaptive systems improve over static. The expert bank is only 0.085 points better overall than one adapter and reduces static-corrected forgetting by 0.156 points (lower is better). Both effects remain below the 0.3-point structural threshold. More importantly, 741 of 742 clips route to expert~0. Thus the observed difference cannot support domain-specialized reuse.

\begin{table*}[t]
\centering
\caption{A--B--C--A revisit results. Corrected forgetting subtracts the static A2--A1 difference; lower is better.}
\label{tab:revisit}
\begin{tabular}{lccccccl}
\toprule
Method & Overall & A1 & A2 & Raw A2--A1 & Corrected & Experts & Routes \\
 & \multicolumn{5}{c}{CER (\%)} & & \\
\midrule
Static & 71.347 & 94.092 & 92.432 & $-1.661$ & 0.000 & 1 & 742 \\
Single adapter & 69.167 & 92.748 & 90.455 & $-2.293$ & $-0.633$ & 1 & 742 \\
Expert bank & 69.081 & 92.748 & 90.299 & $-2.449$ & $-0.789$ & 2 & 741/1 \\
\bottomrule
\end{tabular}
\end{table*}

Post-hoc diagnostics explain the collapse. The route consistency score is 0.9987, which alone looks strong because each labeled domain mostly chooses one expert. Standard clustering purity is only 0.3571 because nearly all domains choose the same expert. Reporting both metrics prevents a collapsed router from being described as domain-pure.

\subsection{Validation-Only Routing Repair Fails the Revisit Gate}

In the pre-fix diagnostic, frozen validation signatures pass the pre-specified proxy gate. The selected order-1 signature has proxy AUROC 0.883, F1 0.859, clustering purity 0.987, and cosine threshold 0.9801. These leave-own-domain-out scores justified one validation revisit run; they are not unseen-domain estimates or results from the strict causal implementation.

\begin{table}[t]
\centering
\caption{Validation A--B--C--A results. Corrected forgetting subtracts the static A2--A1 difference; lower is better.}
\label{tab:validation_revisit}
\begin{tabular}{lrrrr}
\toprule
Method & Overall & A1 & A2 & Corrected \\
 & \multicolumn{4}{c}{CER (\%)} \\
\midrule
Static & 72.329 & 58.505 & 62.713 & 0.000 \\
Single adapter & 69.637 & 57.250 & 56.712 & $-4.746$ \\
Calibrated experts & 71.942 & 58.027 & 62.400 & $+0.165$ \\
\bottomrule
\end{tabular}
\end{table}

The repair nevertheless fails both downstream gates. Relative to one adapter, it worsens overall CER by 2.306 points (95\% CI: 1.700 to 2.917) and worsens the revisit difference-in-differences by 4.911 points (95\% CI: 3.051 to 6.887). It creates all eight allowed experts within A1. Expert~0 receives only 2.34\% of routes, satisfying the anti-collapse criterion, but the A1-dominant expert receives only 42.31\% of A2 routes, below the 80\% reuse criterion. Thus calibration replaces collapse with fragmentation rather than recovering domain-specialized reuse. No repaired-router test run, additional seed, or expert sweep is conducted.

\subsection{Noise Robustness Is Observed but Not Attributed}

With feedback every ten samples, one controlled seed-42 run with 10\% feedback-character replacement changes CER from 66.728\% to 66.770\%, a degradation of 0.041 points. At feedback every 20, the noisy run is 0.023 points better than clean, but this single-seed difference is smaller than the clean three-seed spread and is not interpreted as an improvement. Under that noisy setting, disabling rollback and reliability worsens CER by 0.039 and 0.076 points relative to the full system. These effects are too small to establish robustness or attribute it to either mechanism; repeated corruptions and equivalence intervals remain necessary.

\subsection{Current Decision}

The evidence supports continuing with a single-adapter application paper using full-sequence replay. The error-localized candidate fails its replay gate despite significantly outperforming randomized support. The current dynamic expert implementation is also a no-go under both diagnostic regimes: the original router collapses, while the validation-calibrated router fragments and is significantly worse than one adapter. The pre-specified gates stop localized test evaluation, further expert seeds, and parameter sweeps. A separate next-round protocol decomposes pseudo-only and feedback-only updates and evaluates the full-replay error hybrid on newly locked development and holdout streams. Until that holdout gate is complete, these additions are protocol extensions rather than evidence and do not alter the selected replay result in this section.

\section{Conclusion}
\label{sec:conclusion}

This study evaluated continual personalization for Mandarin VSR under a strict predict-then-update protocol. Across three seeds, no-feedback adaptation of a 75k-parameter residual adapter reduced mean Chinese-LiPS CER by 1.59 absolute points relative to the frozen source model. Supplying one correction every ten utterances added a further 0.832-point mean gain, with the same direction in every matched seed, and produced a 2.42-point total gain over static. These are two valid results under different supervision budgets and should not be collapsed into one claim.

The experiments do not support the current structural implementation. The original dynamic expert bank improved over one adapter by only 0.055 CER points on average and assigned 99.87\% of test-revisit samples to one expert. A validation-only calibration passed its signature proxy gate but then created eight experts, reused the A1-dominant expert on only 42.3\% of A2, and performed significantly worse than one adapter. Because all expert runs could use the known feedback-location flag for pre-prediction growth, they are retained only as pre-fix diagnostics and do not show that dynamic experts are intrinsically ineffective. Reliability and rollback ablations likewise yielded effects too small for a strong causal safety claim.

Validation rejects three superficially plausible update rules. Greedy correct-span preservation is significantly worse than whole-sequence replay, while adapter-only entropy minimization degrades CER by more than 24 points. Target-conditioned error localization significantly outperforms same-mass randomized support, confirming that correction placement carries information, but remains 2.367 CER points worse than replay and also worsens revisit difference-in-differences. It therefore does not enter test evaluation.

On the disjoint target-dev3 development stream, pseudo-only and feedback-only adaptation each significantly improve over static, but adding both update sources does not significantly improve over either source alone. Under identical 75-query budgets, uncertainty selection is 0.272 CER points better than periodic but misses the 0.3-point threshold and has a confidence interval crossing zero; it is 0.021 points worse than random on the point estimate and does not enrich the binary true-error rate. Its static-corrected forgetting interval also fails the pre-registered non-degradation condition. This active-query route is therefore a development no-go, as is the target-dev2 hybrid route. Holdout2 remains frozen and unread, and neither failed route receives extra seeds or tuning.

Before submission, the study still requires stronger matched standard TTA and VSR-personalization baselines, a second real shift type beyond speakers, and a new pre-registered candidate that passes development before any confirmatory holdout run. The resolved gates remove the current expert-growth, error-localized, hybrid, and uncertainty-query implementations from the proposed deployment path. The supported system remains the simpler sparse-feedback single adapter, with its measured resource footprint and limitations; failed candidates are retained as diagnostic evidence rather than promoted through repeated tuning on consumed streams.


\bibliographystyle{IEEEtran}
\bibliography{references}

\end{document}
